\chapter{Stochastic Processes and Markov Chain Monte Carlo}\label{ch:stochproc}

Chapter~\ref{ch:statmech} answered the question \emph{``which distribution?''}
--- Boltzmann's. This chapter answers \emph{``how do we reach it, and how do
we move within it?''}. The answer is the theory of Markov processes: chains
whose future depends on the past only through the present. Markov chains play
a triple role in this thesis: they are the \emph{algorithms} that sample from
intractable distributions (MCMC), they are the \emph{data} whose diffusion
score we compute exactly (the AR(1) sequence introduced in
Section~\ref{sec:ar1}), and --- run in continuous time --- they are the
\emph{corruption and generation processes} of diffusion models
(Chapters~\ref{ch:sde} and~\ref{ch:diffusion}). Standard references for the
material are \citet{robert2004monte} and \citet{brockwell1991time}.

\section{Markov chains: kernels, stationarity, reversibility}
\label{sec:markov}

\begin{definition}[Markov chain]
A sequence of random variables $(X_k)_{k \geq 0}$ with values in a state
space $\mathcal{X}$ is a \emph{Markov chain} if for every $k$ and every
measurable $A \subseteq \mathcal{X}$,
\begin{equation}
  \Prob\big(X_{k+1} \in A \,\big|\, X_k, X_{k-1}, \dots, X_0\big)
  \;=\; \Prob\big(X_{k+1} \in A \,\big|\, X_k\big) .
  \label{eq:markov-property}
\end{equation}
The conditional law is encoded by a \emph{transition kernel}
$M(x' \,|\, x)$: a density (or matrix, for discrete $\mathcal{X}$) with
$\int M(x'|x)\,\dd x' = 1$ for every $x$.
\end{definition}

The Markov property \eqref{eq:markov-property} is exactly a statement about
\emph{factorisation}: the joint law of the first $K$ states is
\begin{equation}
  p(x_0, x_1, \dots, x_{K-1})
  \;=\; p_0(x_0)\, \prod_{k=0}^{K-2} M(x_{k+1} \,|\, x_k) .
  \label{eq:chain-factorisation}
\end{equation}
This one-line factorisation is the single most important equation of the
thesis. It says that the joint distribution of a Markov trajectory is a
product of \emph{local} terms, each touching two consecutive variables ---
the structure that makes the chain a \emph{tree} in the factor-graph sense of
Chapter~\ref{ch:ebm}, that makes its precision matrix tridiagonal in the
Gaussian case of Chapter~\ref{ch:gaussian}, and that makes belief propagation
exact in Chapter~\ref{ch:bp}.

Marginal distributions evolve by integrating the kernel:
$p_{k+1}(x') = \int M(x'|x)\, p_k(x)\,\dd x$, compactly $p_{k+1} = M p_k$.

\begin{definition}[Stationarity]
A distribution $\pi$ is \emph{stationary} (or invariant) for $M$ if
$\pi = M\pi$, i.e.
\begin{equation}
  \pi(x') \;=\; \int M(x' \,|\, x)\, \pi(x)\, \dd x
  \qquad \text{for all } x' .
  \label{eq:stationarity}
\end{equation}
A chain started in $\pi$ stays in $\pi$ forever; the trajectory
$(X_0, X_1, \dots)$ is then a \emph{stationary process}, and all its
finite-dimensional laws are invariant under time shifts.
\end{definition}

\begin{definition}[Detailed balance / reversibility]
The kernel $M$ satisfies \emph{detailed balance} with respect to $\pi$ if
\begin{equation}
  \pi(x)\, M(x' \,|\, x) \;=\; \pi(x')\, M(x \,|\, x')
  \qquad \text{for all } x, x' .
  \label{eq:detailed-balance}
\end{equation}
\end{definition}

Detailed balance says the equilibrium probability flux from $x$ to $x'$
equals the reverse flux --- the chain, watched in equilibrium, is
statistically indistinguishable run forwards or backwards. It is a
\emph{sufficient} condition for stationarity: integrating
\eqref{eq:detailed-balance} over $x$ gives
$\int \pi(x) M(x'|x)\,\dd x = \pi(x') \int M(x|x')\,\dd x = \pi(x')$,
which is \eqref{eq:stationarity}. Its power is practical: detailed balance is
a \emph{local, pointwise} condition, checkable without ever computing an
integral --- which is why every MCMC algorithm in Section~\ref{sec:mcmc} is
engineered to satisfy it.

Stationarity alone does not guarantee that the chain \emph{finds} $\pi$; for
that one needs the chain to be able to reach everywhere
(\emph{irreducibility}) without getting trapped in cycles
(\emph{aperiodicity}). Under these conditions the fundamental convergence
theorem holds \citep{robert2004monte}: the marginal law of $X_k$ converges to
$\pi$ from any initial condition, and time averages converge to expectations
under $\pi$ (the ergodic theorem),
\begin{equation}
  \frac{1}{N} \sum_{k=1}^{N} f(X_k)
  \;\xrightarrow[N \to \infty]{\text{a.s.}}\;
  \E_\pi[f] .
  \label{eq:ergodic}
\end{equation}
Equation \eqref{eq:ergodic} is the rigorous form of the ergodic postulate
invoked in Section~\ref{sec:ensembles}, with the physical dynamics replaced
by an algorithmic one.

\section{The AR(1) chain: the running example of the thesis}\label{sec:ar1}

We now introduce the specific Markov chain that the research chapters solve
completely: the first-order autoregressive (AR(1)) process, the canonical
Gaussian Markov chain \citep{brockwell1991time}. Let
$\alpha \in (-1, 1)$ and let $(\eta_k)$ be i.i.d.\ $\Nd(0, \sigma_\eta^2)$
innovations; define
\begin{equation}
  a_{k+1} \;=\; \alpha\, a_k + \eta_k ,
  \qquad k = 0, 1, \dots, K-2 .
  \label{eq:ar1}
\end{equation}
The transition kernel is Gaussian,
$M(a' | a) = \Nd(a'; \alpha a, \sigma_\eta^2)$, and the joint law of the
trajectory follows from \eqref{eq:chain-factorisation}. Everything about
this chain is computable in closed form; the next toolbox derives the three
facts used throughout.

\begin{toolbox}{Stationary distribution and autocovariance of AR(1)}
\emph{(i) Stationary variance.} Suppose $a_k \sim \Nd(0, \sigma^2)$ and
impose that $a_{k+1}$ has the same variance. Since $\eta_k$ is independent
of $a_k$,
\begin{equation*}
  \operatorname{Var}(a_{k+1})
  = \alpha^2 \operatorname{Var}(a_k) + \sigma_\eta^2
  \;\stackrel{!}{=}\; \operatorname{Var}(a_k)
  \quad\Longrightarrow\quad
  \sigma^2 = \alpha^2 \sigma^2 + \sigma_\eta^2
  \quad\Longrightarrow\quad
  \sigma_\infty^2 = \frac{\sigma_\eta^2}{1 - \alpha^2}.
\end{equation*}
The fixed point exists precisely because $|\alpha| < 1$; the map
$\sigma^2 \mapsto \alpha^2 \sigma^2 + \sigma_\eta^2$ is a contraction, so
\emph{any} initial variance converges to $\sigma_\infty^2$ geometrically.
Since Gaussianity is preserved by the affine recursion, the stationary law is
$\pi = \Nd(0, \sigma_\infty^2)$.

\emph{(ii) Autocovariance.} In the stationary regime, for $d \geq 1$, unroll
the recursion $d$ times:
$a_{k+d} = \alpha^d a_k + \sum_{j=0}^{d-1} \alpha^{j} \eta_{k+d-1-j}$.
All innovation terms are independent of $a_k$, hence
\begin{equation*}
  \operatorname{Cov}(a_k, a_{k+d})
  = \alpha^d\, \operatorname{Var}(a_k)
  = \alpha^d\, \sigma_\infty^2 ,
\end{equation*}
and by symmetry $\operatorname{Cov}(a_i, a_j) = \alpha^{|i-j|}
\sigma_\infty^2$ for all $i, j$: correlations decay \emph{geometrically}
with the lag, with rate $\alpha$. The covariance matrix of $K$ consecutive
frames,
\begin{equation*}
  (\Sigz)_{ij} = \sigma_\infty^2\, \alpha^{|i-j|},
\end{equation*}
depends on $(i,j)$ only through $i - j$: it is a symmetric \emph{Toeplitz}
matrix, a structure we exploit heavily in Chapter~\ref{ch:gaussian}.

\emph{(iii) Detailed balance.} With $\pi = \Nd(0, \sigma_\infty^2)$ and
$M(a'|a) = \Nd(a'; \alpha a, \sigma_\eta^2)$, both sides of
\eqref{eq:detailed-balance} are proportional to
\begin{equation*}
  \exp\!\left[
    -\frac{a^2}{2\sigma_\infty^2}
    -\frac{(a' - \alpha a)^2}{2\sigma_\eta^2}
  \right]
  = \exp\!\left[
    -\frac{(1-\alpha^2) a^2 + (a')^2 - 2\alpha a a' + \alpha^2 a^2}
          {2\sigma_\eta^2}
  \right],
\end{equation*}
and the exponent
$-\big[a^2 + (a')^2 - 2\alpha a a'\big] / (2\sigma_\eta^2)$
is \emph{symmetric} in $(a, a')$. Hence the stationary AR(1) chain is
reversible: statistically, the movie of the chain looks the same played in
either direction. \hfill$\square$
\end{toolbox}

The AR(1) chain is thus the perfect miniature of everything this chapter
discussed: a Markov kernel with an explicit stationary law, geometric
convergence, reversibility, and --- decisive for us --- \emph{exact
Gaussian algebra} at every step. When in Chapter~\ref{ch:gaussian} we corrupt
each frame of this chain with independent noise and ask for the joint score
of the result, every quantity will remain in closed form.

\section{Markov chain Monte Carlo}\label{sec:mcmc}

MCMC inverts the logic of the previous section. There, a kernel $M$ was
given and we asked for its stationary law $\pi$. In sampling problems we are
given $\pi$ --- typically a Boltzmann distribution
$\pi(x) \propto e^{-\beta H(x)}$ known only up to its partition function ---
and we must \emph{design} a kernel $M$ whose stationary law is $\pi$. The
ergodic theorem \eqref{eq:ergodic} then converts a single long trajectory
into expectation values. The idea originates, fittingly, in computational
statistical mechanics \citep{metropolis1953equation}.

\subsection{The Metropolis--Hastings algorithm}

Let $g(x' | x)$ be any proposal kernel we can sample from. The
Metropolis--Hastings (MH) chain \citep{metropolis1953equation,
hastings1970monte} moves from $x$ by (i) proposing $x' \sim g(\cdot | x)$,
(ii) accepting the move with probability
\begin{equation}
  A(x \to x') \;=\; \min\!\left\{ 1,\;
    \frac{\pi(x')\, g(x \,|\, x')}{\pi(x)\, g(x' \,|\, x)}
  \right\},
  \label{eq:mh-acceptance}
\end{equation}
and (iii) staying at $x$ otherwise. Two features make
\eqref{eq:mh-acceptance} remarkable. First, $\pi$ enters only through the
\emph{ratio} $\pi(x')/\pi(x)$: the intractable partition function cancels,
so MH samples from $e^{-\beta H}/Z$ while only ever evaluating energy
differences. Second, the rule is engineered to satisfy detailed balance:

\begin{toolbox}{Detailed balance for Metropolis--Hastings}
The effective transition density for an accepted move $x \neq x'$ is
$M(x'|x) = g(x'|x)\, A(x \to x')$. Compute the flux from $x$ to $x'$:
\begin{equation*}
  \pi(x)\, M(x'|x)
  = \pi(x)\, g(x'|x)\, \min\!\left\{ 1,
    \frac{\pi(x')\, g(x|x')}{\pi(x)\, g(x'|x)} \right\}
  = \min\Big\{ \pi(x)\, g(x'|x),\; \pi(x')\, g(x|x') \Big\},
\end{equation*}
where the second equality uses $c \min\{1, b/c\} = \min\{c, b\}$ for
$c > 0$. The final expression is \emph{symmetric} under
$x \leftrightarrow x'$; therefore
\begin{equation*}
  \pi(x)\, M(x'|x) = \pi(x')\, M(x|x') ,
\end{equation*}
which is detailed balance \eqref{eq:detailed-balance}. The rejection branch
(staying at $x$) contributes only to the diagonal of the kernel and cannot
violate the balance of fluxes between distinct states. Stationarity of $\pi$
follows by the integration argument of Section~\ref{sec:markov}; with an
irreducible, aperiodic proposal, the convergence theorem applies.
\hfill$\square$
\end{toolbox}

\subsection{Gibbs sampling}

When $x = (x_1, \dots, x_n)$ is high-dimensional but each \emph{conditional}
$\pi(x_i \,|\, x_{\setminus i})$ is tractable, the Gibbs sampler
\citep{geman1984stochastic} cycles through the coordinates, resampling each
from its conditional given the current values of the others. Each coordinate
update is an MH move with acceptance probability identically $1$, so detailed
balance holds coordinate-wise and $\pi$ is stationary for the composite
kernel. The Gibbs sampler thrives exactly when the model has sparse
conditional structure --- for a Markov chain,
$\pi(x_k | x_{\setminus k})$ depends only on the neighbours $(x_{k-1},
x_{k+1})$: locality of the model makes the algorithm local. The reader
should hear, already, the theme of Chapter~\ref{ch:bp}: \emph{inference cost
follows graph structure}.

\subsection{Langevin dynamics: MCMC without rejections}

For continuous $x$ and differentiable energy, a different design principle
uses gradients. Consider the discrete-time update
\begin{equation}
  x_{k+1} \;=\; x_k + \epsilon\, \nabla \log \pi(x_k)
  \;+\; \sqrt{2\epsilon}\; \xi_k ,
  \qquad \xi_k \sim \Nd(0, I),
  \label{eq:ula}
\end{equation}
the (unadjusted) \emph{Langevin algorithm}. The drift pushes towards high
probability; the noise maintains exploration. Two observations preview the
next chapters. First, \eqref{eq:ula} only requires the \emph{score}
$\nabla \log \pi$ --- the normalising constant is never needed; this is
precisely why the score is the natural object for generative modelling
(Chapter~\ref{ch:diffusion}), and why so much of this thesis is devoted to
computing scores exactly. Second, as $\epsilon \to 0$ the recursion
\eqref{eq:ula} becomes a stochastic differential equation --- the Langevin
diffusion --- whose stationary law is exactly $\pi$; making that limit
rigorous requires the It\^o calculus and Fokker--Planck theory of
Chapter~\ref{ch:sde}, where we will verify the claim by direct computation.

\section{From algorithmic chains to data chains}

A closing change of perspective, to set up the research problem. In this
chapter Markov chains were \emph{algorithms}: their stationary law was the
target, their trajectory a means to an end. From Chapter~\ref{ch:gaussian}
onwards the chain \eqref{eq:ar1} plays the opposite role: it is the
\emph{data-generating process}, its trajectory $(a_0, \dots, a_{K-1})$ is one
data point --- a ``video'' of $K$ frames --- and the object of interest is
the probability law of the \emph{whole trajectory} after each frame has been
independently corrupted by noise. The factorisation
\eqref{eq:chain-factorisation} then stops being a computational convenience
and becomes the \emph{structure to be discovered}: the entire research
programme of this thesis can be phrased as asking how much of
\eqref{eq:chain-factorisation} survives noising, and how an inference
algorithm --- or a neural network --- can exploit what survives.
